% GENERATED FILE. Edit canonical lessons, then run npm run generate:books.
% canonical-source-hash: fnv1a64:38406cf9057ea67c
% canonical-lessons: UR-C395-tayar, UR-C395-pankchar, UR-C395-takkar, UR-C395-helmet, UR-C395-pichakna

\chapter{Tyres and Punctures}
\label{ch:ur-tyres-and-punctures}
\hlchaptermodality{395}

\begin{chapteropening}
\textbf{By the end of this chapter:} \emph{I can deal with a flat tyre, a knock and a dent, and a helmet for the ride.}

Complete the last lesson of chapter 395: I can deal with a flat tyre, a knock and a dent, and a helmet for the ride.
\end{chapteropening}

\section[\texorpdfstring{\ur{ٹائر} (ṭāyar)}{ٹائر (ṭāyar)}]{\ur{ٹائر} (ṭāyar) — a tyre}

\label{lesson:UR-C395-tayar}



\begin{quote}
Before the new one: say the Urdu for to bounce, then the Urdu for deserted.
\end{quote}

\subsection*{You'll want to know: \ur{ٹائر}}
\textbf{\ur{ٹائر}} — \emph{ṭāyar} — \textquotedblleft{}a tyre\textquotedblright{}.

\begin{quote}
\ur{گاڑی} \ur{کا} \ur{ایک} \ur{ٹائر} \ur{پرانا} \ur{ہے۔}
\end{quote}

\begin{quote}
\emph{gāṛī kā ek ṭāyar purānā hai.}
\end{quote}

\begin{quote}
One of the car's tyres is old.
\end{quote}

Masculine.

\begin{cousinweb}
\textbf{\ur{ٹائر}} (\emph{ṭāyar}) is borrowed from English \emph{tyre}.
\end{cousinweb}

\begin{grammarlens}[title={What you've built}]
One more word for getting around.
\end{grammarlens}

\subsection*{Your turn}
\begin{itemize}
\raggedright
  \item \emph{Say it:} \emph{ṭāyar}
  \item \emph{Say it:} \emph{ṭāyar}, once more
  \item \emph{Say it:} \emph{gāṛī kā ek ṭāyar purānā hai.} — One of the car's tyres is old.
  \item \emph{Recall:} say the Urdu for to bounce, then the Urdu for deserted, then say \emph{ṭāyar} again
\end{itemize}

\begin{tcolorbox}[breakable,colback=teal!4,colframe=teal!35!black,title={Before you move on}]
What does \ur{ٹائر} mean? (\textquotedblleft{}A tyre\textquotedblright{}.) Say it once more.
\end{tcolorbox}
\section[\texorpdfstring{\ur{پنکچر} (pankchar)}{پنکچر (pankchar)}]{\ur{پنکچر} (pankchar) — a puncture}

\label{lesson:UR-C395-pankchar}



\begin{quote}
Before the new one: say the Urdu for deserted, then the Urdu for a tyre.
\end{quote}

\subsection*{You'll want to know: \ur{پنکچر}}
\textbf{\ur{پنکچر}} — \emph{pankchar} — \textquotedblleft{}a puncture\textquotedblright{}.

\begin{quote}
\ur{راستے} \ur{میں} \ur{ہماری} \ur{گاڑی} \ur{کا} \ur{ٹائر} \ur{پنکچر} \ur{ہو} \ur{گیا۔}
\end{quote}

\begin{quote}
\emph{rāste meṅ hamārī gāṛī kā ṭāyar pankchar ho gayā.}
\end{quote}

\begin{quote}
On the way our car got a flat tyre.
\end{quote}

Masculine.

\begin{cousinweb}
\textbf{\ur{پنکچر}} (\emph{pankchar}) is borrowed from English \emph{puncture}; with \textbf{\ur{ہونا}} it means a tyre goes flat.
\end{cousinweb}

\begin{grammarlens}[title={What you've built}]
One more word for getting around.
\end{grammarlens}

\subsection*{Your turn}
\begin{itemize}
\raggedright
  \item \emph{Say it:} \emph{pankchar}
  \item \emph{Say it:} \emph{pankchar}, once more
  \item \emph{Say it:} \emph{rāste meṅ hamārī gāṛī kā ṭāyar pankchar ho gayā.} — On the way our car got a flat tyre.
  \item \emph{Recall:} say the Urdu for deserted, then the Urdu for a tyre, then say \emph{pankchar} again
\end{itemize}

\begin{tcolorbox}[breakable,colback=teal!4,colframe=teal!35!black,title={Before you move on}]
What does \ur{پنکچر} mean? (\textquotedblleft{}A puncture\textquotedblright{}.) Say it once more.
\end{tcolorbox}
\section[\texorpdfstring{\ur{ٹکر} (ṭakkar)}{ٹکر (ṭakkar)}]{\ur{ٹکر} (ṭakkar) — a collision; a knock}

\label{lesson:UR-C395-takkar}



\begin{quote}
Before the new one: say the Urdu for a tyre, then the Urdu for a puncture.
\end{quote}

\subsection*{You'll want to know: \ur{ٹکر}}
\textbf{\ur{ٹکر}} — \emph{ṭakkar} — \textquotedblleft{}a collision; a knock\textquotedblright{}.

\begin{quote}
\ur{پل} \ur{پر} \ur{دو} \ur{گاڑیوں} \ur{کی} \ur{ٹکر} \ur{ہوئی۔}
\end{quote}

\begin{quote}
\emph{pul par do gāṛiyoṅ kī ṭakkar huī.}
\end{quote}

\begin{quote}
Two cars collided on the bridge.
\end{quote}

Feminine.

\begin{cousinweb}
\textbf{\ur{ٹکر}} (\emph{ṭakkar}) is inherited: CDIAL lists Hindi \emph{ṭakkar}, \textquotedblleft{}a blow\textquotedblright{}, under \emph{ṭakkarā}, \textquotedblleft{}a blow on the head\textquotedblright{} (CDIAL 5424).
\end{cousinweb}

\begin{grammarlens}[title={What you've built}]
One more word for getting around.
\end{grammarlens}

\subsection*{Your turn}
\begin{itemize}
\raggedright
  \item \emph{Say it:} \emph{ṭakkar}
  \item \emph{Say it:} \emph{ṭakkar}, once more
  \item \emph{Say it:} \emph{pul par do gāṛiyoṅ kī ṭakkar huī.} — Two cars collided on the bridge.
  \item \emph{Recall:} say the Urdu for a tyre, then the Urdu for a puncture, then say \emph{ṭakkar} again
\end{itemize}

\begin{tcolorbox}[breakable,colback=teal!4,colframe=teal!35!black,title={Before you move on}]
What does \ur{ٹکر} mean? (\textquotedblleft{}A collision; a knock\textquotedblright{}.) Say it once more.
\end{tcolorbox}
\section[\texorpdfstring{\ur{ہیلمٹ} (helmeṭ)}{ہیلمٹ (helmeṭ)}]{\ur{ہیلمٹ} (helmeṭ) — a helmet}

\label{lesson:UR-C395-helmet}



\begin{quote}
Before the new one: say the Urdu for a puncture, then the Urdu for a collision.
\end{quote}

\subsection*{You'll want to know: \ur{ہیلمٹ}}
\textbf{\ur{ہیلمٹ}} — \emph{helmeṭ} — \textquotedblleft{}a helmet\textquotedblright{}.

\begin{quote}
\ur{موٹر} \ur{سائیکل} \ur{پر} \ur{ہیلمٹ} \ur{پہننا} \ur{چاہیے۔}
\end{quote}

\begin{quote}
\emph{moṭar sāikil par helmeṭ pahannā chāhiye.}
\end{quote}

\begin{quote}
You should wear a helmet on a motorbike.
\end{quote}

Masculine.

\begin{cousinweb}
\textbf{\ur{ہیلمٹ}} (\emph{helmeṭ}) is borrowed from English \emph{helmet}.
\end{cousinweb}

\begin{grammarlens}[title={What you've built}]
One more word for getting around.
\end{grammarlens}

\subsection*{Your turn}
\begin{itemize}
\raggedright
  \item \emph{Say it:} \emph{helmeṭ}
  \item \emph{Say it:} \emph{helmeṭ}, once more
  \item \emph{Say it:} \emph{moṭar sāikil par helmeṭ pahannā chāhiye.} — You should wear a helmet on a motorbike.
  \item \emph{Recall:} say the Urdu for a puncture, then the Urdu for a collision, then say \emph{helmeṭ} again
\end{itemize}

\begin{tcolorbox}[breakable,colback=teal!4,colframe=teal!35!black,title={Before you move on}]
What does \ur{ہیلمٹ} mean? (\textquotedblleft{}A helmet\textquotedblright{}.) Say it once more.
\end{tcolorbox}
\section[\texorpdfstring{\ur{پچکنا} (pichaknā)}{پچکنا (pichaknā)}]{\ur{پچکنا} (pichaknā) — to get dented; to be squashed}

\label{lesson:UR-C395-pichakna}



\begin{quote}
Before the new one: say the Urdu for a collision, then the Urdu for a helmet.
\end{quote}

\subsection*{You'll want to know: \ur{پچکنا}}
\textbf{\ur{پچکنا}} — \emph{pichaknā} — \textquotedblleft{}to get dented; to be squashed\textquotedblright{}.

\begin{quote}
\ur{ٹکر} \ur{میں} \ur{گاڑی} \ur{کا} \ur{دروازہ} \ur{پچک} \ur{گیا۔}
\end{quote}

\begin{quote}
\emph{ṭakkar meṅ gāṛī kā darvāzā pichak gayā.}
\end{quote}

\begin{quote}
The car door got dented in the collision.
\end{quote}

Past (it …): \ur{وہ} \ur{ڈبہ} \ur{پچک} \ur{گیا}, \emph{voh ḍabbā pichak gayā}.

\begin{cousinweb}
\textbf{\ur{پچکنا}} (\emph{pichaknā}) is inherited: CDIAL lists the Hindi cousin \emph{picaknā}, \textquotedblleft{}to be pressed\textquotedblright{}, under Sanskrit \emph{piccayati}, \textquotedblleft{}presses flat\textquotedblright{} (CDIAL 8149); the sense \textquotedblleft{}to get dented\textquotedblright{} is probably a later one.
\end{cousinweb}

\begin{grammarlens}[title={What you've built}]
One more word for getting around.
\end{grammarlens}

\subsection*{Your turn}
\begin{itemize}
\raggedright
  \item \emph{Say it:} \emph{pichaknā}
  \item \emph{Say it:} \emph{pichaknā}, once more
  \item \emph{Say it:} \emph{ṭakkar meṅ gāṛī kā darvāzā pichak gayā.} — The car door got dented in the collision.
  \item \emph{Recall:} say the Urdu for a collision, then the Urdu for a helmet, then say \emph{pichaknā} again
\end{itemize}

\begin{tcolorbox}[breakable,colback=teal!4,colframe=teal!35!black,title={Before you move on}]
What does \ur{پچکنا} mean? (\textquotedblleft{}To get dented; to be squashed\textquotedblright{}.) Say it once more.
\end{tcolorbox}
